\documentclass[11pt]{article}
\usepackage[a4paper,margin=23mm]{geometry}
\usepackage{amsmath,amssymb,xcolor,hyperref,fancyhdr}
\usepackage{fontspec}
\setmainfont{Times New Roman}
\setmonofont{Times New Roman}
\usepackage[italic]{mathastext}
\setlength{\emergencystretch}{3em}
\setlength{\parskip}{7pt}\setlength{\parindent}{0pt}
\pagestyle{fancy}\fancyhf{}\setlength{\headheight}{14pt}
\lhead{FlowMLLab}\rhead{Reading and experiment guide}\cfoot{\thepage}
\begin{document}
\noindent\textbf{Week 10 | Reading the evidence}\par\medskip\hrule\medskip
\section*{Trace every figure to the run that produced it}
Retained author figures document a fixed experiment. Newly regenerated figures
must be displayed from the fresh output directory, next to that run's metrics.
The default cavity reproduction interpolates fields in log Knudsen number; it
does not retrain the article's per-Knudsen MLPs.

For two training Knudsen numbers $Kn_a,Kn_b$, the interpolation weight is

\[w=\frac{\log Kn-\log Kn_a}{\log Kn_b-\log Kn_a},\qquad
q=(1-w)q_a+wq_b.\]

$q$ is a field and $w$ is dimensionless. This is interpolation for $0\le w\le1$;
outside that interval it is extrapolation. Verify the weight by substitution.
For error $\|q-\hat q\|_2/S$, state the scale $S$. A velocity-component error
normalized by the lid speed can look small while its component-relative error
is large. Neither normalization may be silently substituted for the other.
Also inspect heat flux, stress and translational-temperature overshoot: matching
density alone does not validate nonequilibrium shock structure.
\paragraph{Check your understanding.}
Identify the input information, the reference, the metric denominator and the
evaluation split in one worked notebook comparison. Predict which conclusion
would change if any one of those four changed, then check your prediction.
\end{document}
